% !TeX root = ../main.tex
% ---------------------------------------------------------------------------
%  Sustainability and Carbon in Underground Construction
%  Enriched from the author's study report on sustainability and resilience of
%  underground infrastructure (Lecture Slides/
%  sustainability_resilience_logic_study_report.tex), which is the primary
%  source for the module-based carbon workflow and the design-variable
%  hierarchy below.
%  Other sources: EN 15978; EN 17472; PAS 2080; ITA WG on sustainability;
%  Infrastructure Sustainability Council (Australia).
% ---------------------------------------------------------------------------
\chapter{Sustainability and Carbon in Underground Construction}
\label{ch:sustainability}

Underground works are material-intensive and energy-intensive. A bored tunnel
consumes concrete, steel, grout and backfill at a rate few other structures
match, and consumes it in a setting where the excavation energy itself is a
significant emission source. Carbon is therefore not a reporting exercise
appended to a completed design: it is determined by decisions about alignment,
diameter, lining thickness, support timing and excavation method --- the same
decisions that govern safety and serviceability.

This chapter treats carbon accordingly. It sets out the life-cycle module
framework that makes a carbon figure defensible, quantifies the
construction-stage terms that dominate a tunnel, and then treats carbon as a
design variable alongside factor of safety and deformation.

\section{Sustainability drivers and carbon sources}
\label{sec:sustainability-drivers-and-carbon-sources}

\subsection{Drivers for underground works}
\label{sec:sustainability-drivers}
\scopenote{Regulatory, client and social drivers; net zero commitments; and
the position of underground infrastructure in decarbonising transport and
energy systems.}

\subsection{Where the carbon is}
\label{sec:embodied-carbon}
\scopenote{Concrete and cement above all, then steel reinforcement, then spoil
haulage and construction energy. Typical proportions by tunnel type, and why
the ranking shifts between a segmentally lined TBM drive and a sprayed
concrete lined cavern.}

\section{Life-cycle carbon assessment}
\label{sec:lifecycle-carbon}

\subsection{Module boundaries}
\label{sec:carbon-modules}
\scopenote{The EN 15978 and EN 17472 module structure applied to tunnels:
A1--A3 product, A4 transport, A5 construction, B use and maintenance, C end of
life, D beyond-system benefits. Most published tunnel carbon figures are not
comparable because the boundary is not stated.}

\subsection{Model-based quantity take-off}
\label{sec:model-based-carbon}
\scopenote{Carbon assessment depends on a quantity model, a carbon-factor
library, a declared system boundary and a construction process model. Taking
quantities from the digital asset model of
Chapter~\ref{ch:digital-underground} --- lining concrete, reinforcement,
grout, backfill, fit-out, spoil --- rather than from spreadsheets rebuilt at
each design stage.}

\subsection{Carbon factors and their provenance}
\label{sec:carbon-factors}
\scopenote{Database factors, environmental product declarations and
project-specific data. Concrete grade, supplementary cementitious material
content, steel type and geographical context must be recorded with the factor,
or the result cannot be audited.}

\section{Construction-stage carbon}
\label{sec:construction-stage-carbon}

\subsection{Accounting logic}
\label{sec:carbon-formula}
The construction-stage total assembles a product term, a transport term and a
process term over the material inventory:
\[
GWP_A \;=\; \sum_i V_i \left( GWP_{i,\mathrm{A1-A3}} + GWP_{i,\mathrm{A4}} \right)
\;+\; GWP_{\mathrm{A5}}
\]
where $V_i$ is the quantity of material $i$.

\begin{warnbox}[This is accounting logic, not a design equation]
The expression above is a bookkeeping structure, not a closed-form predictor.
Its value lies in forcing every reported figure to be traceable to a quantity,
a boundary, a carbon factor and a stated modelling assumption. A carbon number
quoted without those four things cannot be compared with another project's.
\end{warnbox}

\subsection{Transport and spoil haulage}
\label{sec:carbon-transport}
\scopenote{A4 emissions from haulage distance, mode, mass, trip count and
vehicle factor. Spoil haulage is frequently the largest single A4 term, and it
is set by disposal siting decisions taken outside the engineering team.}

\subsection{Excavation energy}
\label{sec:carbon-tbm-energy}
For mechanised excavation the construction-stage term is dominated by machine
energy, which can be tied to the electricity and fuel factors applying to the
site supply:
\[
GWP_{\mathrm{A5}} \;=\; E_{\mathrm{TBM}} \left( EF_2 + EF_3 \right)
\]
\scopenote{$E_{\mathrm{TBM}}$ should be estimated from ground--machine
interaction --- torque, thrust, rate of penetration, rotation speed, diameter,
depth, ground class and in-situ stress --- rather than assigned as a generic
percentage of project emissions. This links directly to the performance
prediction methods of Chapter~\ref{ch:tbm-engineering}.}

\section{Carbon as a design variable}
\label{sec:carbon-as-design-variable}

\subsection{The design-variable hierarchy}
\label{sec:carbon-hierarchy}
Carbon sensitivity is not distributed evenly across design decisions. In
approximate order of leverage:

\begin{enumerate}[nosep]
  \item alignment, depth, diameter and geometric scale;
  \item lining thickness, concrete grade, SCM content and reinforcement ratio;
  \item support distance from the face, and the ground--support interaction it
        mobilises;
  \item excavation method and the disturbance it causes;
  \item machine performance --- thrust, torque, rate of penetration and the
        resulting construction energy;
  \item spoil removal distance and transport mode;
  \item restoration strategy and its timing during operation.
\end{enumerate}

\scopenote{The practical message is that early geometric decisions and
construction control usually dominate the achievable reduction, and material
substitution alone is a late and limited lever.}

\subsection{Geo-structural carbon optimisation}
\label{sec:carbon-optimisation}
\scopenote{A defensible optimisation workflow: define design variables; define
geotechnical inputs and their uncertainty ranges; define performance
constraints on factor of safety, support pressure, deformation, settlement,
constructability and durability; define the carbon objective by module; then
evaluate alternatives using analytical, convergence--confinement, numerical or
surrogate models according to project stage.}

\subsection{Trade-offs and Pareto solutions}
\label{sec:carbon-pareto}
\scopenote{There is rarely a single minimum-carbon design. Multi-objective
optimisation yields a Pareto set, from which the preferred option is the one
satisfying safety and deformation constraints without unnecessary support
demand. Allowing controlled ground deformation can reduce structural demand
and carbon together --- but only within serviceability and construction-risk
limits. Applied in Chapter~\ref{ch:integrated-design}.}

\section{Materials, spoil and circularity}
\label{sec:materials-and-spoil}

\subsection{Low-carbon concrete}
\label{sec:low-carbon-concrete}
\scopenote{Supplementary cementitious materials, alkali-activated and
geopolymer binders, and the tension between carbon reduction and the early
strength that tunnelling cycles depend on
(Section~\ref{sec:early-age}).}

\subsection{Spoil management, reuse and circularity}
\label{sec:spoil}
\scopenote{Excavated material as a resource rather than waste: aggregate
recovery, conditioning residues, contamination constraints and the regulatory
classification of tunnel spoil.}

\section{Whole-life and thermal energy}
\label{sec:whole-life-and-thermal-energy}

\subsection{Operational energy and whole-life carbon}
\label{sec:operational-carbon}
\scopenote{Ventilation, lighting, pumping and maintenance over a century of
operation, and the trade-off against embodied carbon in decisions such as
drained versus tanked linings.}

\subsection{Energy geostructures}
\label{sec:energy-geostructures}
\scopenote{Thermally activated linings and absorber piles, geothermal recovery
from tunnels, and the design implications of cyclic thermal loading.}

\section{Urban context, approvals and reporting}
\label{sec:urban-and-environmental-context}

\subsection{Underground space and urban sustainability}
\label{sec:underground-space}
\scopenote{Land take, surface amenity, urban heat and resilience: the case for
underground solutions on grounds other than cost.}

\subsection{Environmental impact and approvals}
\label{sec:environmental-impact}
\scopenote{Groundwater, noise, vibration, spoil transport and ecological
impacts within the environmental assessment frameworks introduced in
Chapter~\ref{ch:context}.}

\subsection{Rating schemes and reporting}
\label{sec:rating-schemes}
\scopenote{Infrastructure Sustainability Council ratings and comparable
international schemes; what they measure and what they miss.}
